\documentclass{article}
\begin{document}
\section*{Compatibility and scope: final-grade decomposition}
The binary hom comparison of PR 182 preserves the independently defined identity square and pasting of arbitrary squares over composable base arrows. The equalities are in the binary apex hom-set: proving only projected equalities would leave the transported second grade unchecked. The proof handles that transport by dependent-pair extensionality and heterogeneous equality of the grade components.
\par Cutting off the last grade is inverse to right append for every nonempty independent chain. Squares between two appended chains are equivalent to a prefix square over the specified base arrow and a final grade arrow over its computed terminal map. The outgoing map agrees with successor transport applied to this last grade arrow. No endpoint splitting is used by these decompositions; no grade arrow is required to be invertible.
\par With the existing split endpoint laws the recursively presented positive power has an object equivalence with all independent chains of the corresponding length, preserving both endpoint objects. Length index zero in the positive tower means one primitive span. The successor transports the actual grade along the established prefix endpoint equality; neither an equipment associator nor an equivalence of span powers is assumed.
\par Missing: induction on native tower arrows comparing the endpoint transports, general identity and composition preservation, and the resulting functorial equivalence. The prior strict universal property applies to the equality-matched model; its comparison with the presented apex is an equivalence, not a claimed strict inverse. No empty-span or monad unit law, flattening, strict-algebra recovery or closure of issue 105 follows.
\end{document}
